\documentclass[10pt,a4paper]{amsart}
\usepackage[margin=19mm]{geometry}
\usepackage{amsmath,amssymb,mathtools}
\usepackage[scr=rsfs,cal=euler]{mathalfa}
\usepackage{xfrac}
\usepackage[unicode,hidelinks,bookmarks=false]{hyperref}
\newcommand{\ie}{i.e.\ }
\newcommand{\cf}{cf.\ }
\newcommand{\resp}{respectively\ }
\newcommand{\Subsection}{\subsection}
\providecommand{\nobreakdash}{\nobreak\hbox{-}}
\setlength{\emergencystretch}{2em}
\multlinegap0pt
\usepackage{xcolor}
\usepackage{amssymb}
\newtheorem{lemma}{Lemma}
\newcommand{\C}{\mathbb{C}}
\newcommand{\D}{\mathbb{D}}
\renewcommand{\Re}{\operatorname{Re}}
\newcommand{\T}{\mathbb{T}}
\newcommand{\Z}{\mathbb{Z}}
\title{The norm of the backward shift on $H^4$ is $\sqrt[4]{\varphi}$}
\author{Konstantinos Bampouras, Adri\'an Llinares}
\date{Source: arXiv:2605.10469v1 (CC BY 4.0)}
\begin{document}
\maketitle

\noindent\emph{Source and licence.} The norm of the backward shift on $H^4$ is $\sqrt[4]{\varphi}$. Version arXiv:2605.10469v1 (CC BY 4.0), distributed by arXiv under the Creative Commons Attribution 4.0 International licence, http://creativecommons.org/licenses/by/4.0/, as recorded in the arXiv licence metadata for this exact version and shown on the abstract page https://arxiv.org/abs/2605.10469v1. Full licence notice: \texttt{licenses/arxiv-2605.10469-CC-BY-4.0.txt}.\par\smallskip

\noindent\textit{Editorial scope.} Counted unit: the article's lemma lemma (source0.tex lines 122--128). The article's complete original proof of it is delivered with it (source0.tex lines 130--175, one \texttt{proof} environment of 46 lines). No mathematics is written, repaired or re-derived: the statement and the proof are the article's own lines, in the article's own notation, and no retained line is substituted or re-flowed.  No further source-local context is needed: every cross-reference of every retained line resolves inside the two delivered spans.  Nothing was omitted from any retained span, so the package carries no omission record.  No retained line contains a citation, so the wrapper carries no bibliography and no bibliography entry.  No retained line contains a figure environment, an image-inclusion command or drawing code, so no image is delivered and the package declares no asset dependency.  Editorial formatting only, in addition: an AMS \texttt{amsart} preamble that restates the article's own declarations and macros used by the retained lines, namely the environments \texttt{lemma} on separate counters, together with the standard packages those lines need.  The retained environments print in this document as Lemma 1 in order of appearance, whereas the article prints the same items with the numbering of its own full document.  The complete native source of the article as distributed in the arXiv e-print for this exact version is bundled unchanged as \texttt{sources/200297/source0.tex}.
        \begin{lemma}\label{lemma}
            For every $f \in H^4$ we have that
            \begin{equation} \label{eqn:CubicBound}
                |\langle f^2, Bf \rangle|^2 \leq \dfrac{1}{2} \| f \|_{H^2}^2 ( \| f \|_{H^4}^4-\|f\|_{H^2}^4 ),
            \end{equation}
            where $\langle \cdot, \cdot \rangle$ denotes the inner product of $L^2 (\T)$. Moreover, equality holds if and only if $f$ is a constant multiple of the backward shift of an inner function.
        \end{lemma}

        \begin{proof}
            
            For almost every point of $\T$ we have that
            \[
            Bf(z) = \overline{z} \big( f(z) - f(0) \big).
            \]
            Since $z f^2(z)$ vanishes at the origin, we have that
            \begin{align*}
                \langle f^2, Bf \rangle & = \dfrac{1}{2\pi} \int_0^{2\pi} f^2 (e^{it}) \overline{Bf(e^{it})} \, dt \\
                & = \dfrac{1}{2\pi} \int_0^{2\pi} e^{it} f^2 (e^{it}) \big( \overline{f(e^{it})} - \overline{f(0)} \big) \, dt \\
                & = \dfrac{1}{2\pi} \int_{0}^{2\pi} e^{it} |f(e^{it})|^2 f(e^{it}) \, dt.
            \end{align*}
            If $P_+$ represents the Riesz projection (i.e., the orthogonal projection from $L^2(\T)$ onto $H^2$) we will have that the following identity holds:
            \[
            \langle f^2, B f \rangle = \langle f, P_+ ( \overline{z} |f|^2 ) \rangle.
            \]
            Hence, Cauchy-Schwarz inequality yields that
            \begin{equation} \label{eqn:CS}
                | \langle f^2, B f \rangle | \leq \| f \|_{H^2} \| P_+ ( \overline{z} |f|^2 ) \|_{H^2}.
            \end{equation}
            
            Let $\{ b_k \}_{k \in \Z}$ be the Fourier coefficients of $|f|^2 \in L^2(\T)$, which is a real-valued function and therefore we have that $b_k = \overline{b_{-k}}$ for all $k < 0$. Observe that $\| f \|_{H^2}^2 = b_0$ and
            \[
            \| P_+ (\overline{z}|f|^2) \|_{H^2}^2 = \sum_{k = 1}^\infty |b_k|^2 = \dfrac{\|f\|_{H^4}^4 - \| f \|_{H^2}^4}{2},
            \]
            which proves that inequality \eqref{eqn:CubicBound} holds.
            
            Now, assume that equality in \eqref{eqn:CubicBound} is attained for certain function $f \in H^4$. In particular, equality is also attained in \eqref{eqn:CS} and therefore there exists a constant $\lambda \in \C$ such that $P_+ (\overline{z}|f|^2) = \lambda f$.
            
            If $g(z) = z f(z)$, the identity above becomes
            \[
            \lambda g(z) = \sum_{k = 1}^\infty b_k z^k
            \]
            and then
            \[
            |g(z)|^2 = |f(z)|^2 = \|f\|_{H^2}^2 + 2 \Re \{ \lambda g(z) \}
            \]
            almost everywhere in $\T$. In other words, the function $g - \overline{\lambda} \in H^4$ has constant modulus almost everywhere and that means there exists an inner function $I$ and $\mu \in \C$ such that
            \[
            g(z) - \overline{\lambda} = \mu I(z), \quad z \in \D.
            \]
            Finally, $g$ vanishes at 0 so
            \[
            f = Bg = \mu B I.
            \]
        \end{proof}

\end{document}
